\documentclass[
10pt, % Main document font size
a4paper, % Paper type, use 'letterpaper' for US Letter paper
oneside, % One page layout (no page indentation)
%twoside, % Two page layout (page indentation for binding and different headers)
headinclude,footinclude, % Extra spacing for the header and footer
BCOR5mm, % Binding correction
]{scrartcl}
\input{stru/structure.tex}
\hyphenation{Fortran hy-phen-ation} 

\title{\normalfont{FTML practical session 13 solution}} 
\author{\spacedlowsmallcaps{}
} 
\date{\today}
\begin{document}
\renewcommand{\sectionmark}[1]{\markright{\spacedlowsmallcaps{#1}}}
\lehead{\mbox{\llap{\small\thepage\kern1em\color{halfgray} \vline}\color{halfgray}\hspace{0.5em}\rightmark\hfil}} 
\pagestyle{scrheadings}
\maketitle 
\setcounter{tocdepth}{3} 

\section{\large\color{Blue}Heavy-ball}

\exercice{}

We have that 

\begin{equation}
    \begin{aligned}
        \label{eq:}
	\theta_{t+1}- \theta^* &= \theta_{t}- \theta^*- \gamma \nabla_{ \theta}f(\theta_t)+ \beta( \theta_{t}- \theta_{t-1})  \\
	&= \theta_{t}- \theta^*- \gamma ( H \theta_{t}- \frac{1}{n} X^T y)+ \beta( \theta_{t}- \theta_{t-1})  \\
	&= \theta_{t}- \theta^*- \gamma ( H \theta_{t}- H \theta^*)+ \beta( \theta_{t}- \theta_{t-1})  \\
	&= (I_d- \gamma H)(\theta_{t}- \theta^*)+ \beta\Big( (\theta_{t}- \theta^*)- (\theta_{t-1}-\theta^*)\Big)  \\
    \end{aligned}
\end{equation}

Hence,

\begin{equation}
    \begin{aligned}
        \label{eq:}
	\langle \theta_{t+1}- \theta^*, u_{\lambda} \rangle &=     \langle (I_d- \gamma H)(\theta_{t}- \theta^*)+ \beta\Big( (\theta_{t}- \theta^*)- (\theta_{t-1}-\theta^*)\Big), u_{\lambda} \rangle\\
	&= \langle (I_d- \gamma H)(\theta_{t}- \theta^*), u_{\lambda} \rangle+ \beta\langle \theta_{t}- \theta^*, u_{\lambda} \rangle- \beta\langle\theta_{t-1}-\theta^*, u_{\lambda} \rangle\\
	&= \langle \theta_{t}- \theta^*, (I_d- \gamma H)u_{\lambda} \rangle+ \beta\langle \theta_{t}- \theta^*, u_{\lambda} \rangle- \beta\langle\theta_{t-1}-\theta^*, u_{\lambda} \rangle\\
	&= \langle \theta_{t}- \theta^*, (1- \gamma \lambda)u_{\lambda} \rangle+ \beta\langle \theta_{t}- \theta^*, u_{\lambda} \rangle- \beta\langle\theta_{t-1}-\theta^*, u_{\lambda} \rangle\\
	&= (1- \gamma \lambda+ \beta)\langle \theta_{t}- \theta^*, u_{\lambda} \rangle - \beta\langle\theta_{t-1}-\theta^*, u_{\lambda} \rangle\\
    \end{aligned}
\end{equation}

which is the expected result. We have used that $ I_d- \gamma H$ is symmetric and that $ u_{\lambda}$ is an eigenvector of $H$.
\\

\exercice{}
The characteristic polynomial is

\begin{equation}
    P(X) = X^2-(1- \gamma\lambda +\beta)X+ \beta
\end{equation}

In order to find the values of the sequence $(a_t)_{t\in \mathbb{R} n}$ we need to know the roots of the equation

\begin{equation}
    \label{eq:cara}
    X^2-(1- \gamma\lambda +\beta)X+ \beta=0
\end{equation}

The discriminant writes

\begin{equation}
    \begin{aligned}
        \label{eq:}
	\Delta &= (1- \gamma\lambda +\beta)^2- 4\beta \\
	&= 1+ (\gamma\lambda)^2 +\beta^2-2 \gamma\lambda + 2 \beta - 2 \gamma\lambda \beta- 4\beta \\
	&= 1+ (\gamma\lambda)^2 +\beta^2-2 \gamma\lambda  - 2 \gamma\lambda \beta- 2\beta \\
	&= 1+ (\frac{4}{( \sqrt{L} + \sqrt{\mu} )^2}\lambda)^2 +\Big( \frac{\sqrt{L} - \sqrt{\mu}}{\sqrt{L} + \sqrt{\mu}}  \Big)^4\\
	&-2 \frac{4}{( \sqrt{L} + \sqrt{\mu} )^2}\lambda - 2 \frac{4}{( \sqrt{L} + \sqrt{\mu} )^2}\lambda \Big( \frac{\sqrt{L} - \sqrt{\mu}}{\sqrt{L} + \sqrt{\mu}}  \Big)^2- 2\Big( \frac{\sqrt{L} - \sqrt{\mu}}{\sqrt{L} + \sqrt{\mu}}  \Big)^2 \\
    \end{aligned}
\end{equation}

If we use the common denominator $( \sqrt{L} + \sqrt{\mu} )^4$, the numerator $N$ writes

\begin{equation}
    \begin{aligned}
        \label{eq:}
	N&= ( \sqrt{L} + \sqrt{\mu} )^4+16 \lambda^2+( \sqrt{L} - \sqrt{\mu} )^4\\
	&-8\lambda( \sqrt{L} + \sqrt{\mu} )^2-8\lambda ( \sqrt{L} - \sqrt{\mu} )^2-  2 ( \sqrt{L} + \sqrt{\mu} )^2( \sqrt{L} - \sqrt{\mu} )^2\\
    \end{aligned}
\end{equation}

But

\begin{equation}
    \begin{aligned}
        \label{eq:}
	( \sqrt{L} + \sqrt{\mu} )^4+( \sqrt{L} - \sqrt{\mu} )^4-  2 ( \sqrt{L} + \sqrt{\mu} )^2( \sqrt{L} - \sqrt{\mu} )^2&= \Big(( \sqrt{L} + \sqrt{\mu} )^2- ( \sqrt{L} - \sqrt{\mu} )^2 \Big)^2\\
	&= \Big( (L+\mu+2 \sqrt{L\mu})-(L+\mu-2 \sqrt{L\mu})  \Big)^2\\
	&= \Big( 4 \sqrt{L\mu}   \Big)^2\\
	&= 16L\mu
    \end{aligned}
\end{equation}

And 

\begin{equation}
   \begin{aligned}
       \label{eq:}
       -8\lambda( \sqrt{L} + \sqrt{\mu} )^2-8\lambda ( \sqrt{L} - \sqrt{\mu} )^2&= -8\lambda\Big(( \sqrt{L} + \sqrt{\mu} )^2+( \sqrt{L} - \sqrt{\mu} )^2\Big)\\
       &= -8\lambda\Big(2 (L+\mu)\Big)\\
       &= -16\lambda(L+\mu)\\
   \end{aligned} 
\end{equation}

Finally, 

\begin{equation}
    \begin{aligned}
        \label{eq:}
	N &= 16\lambda^2+16\lambda \mu -16\lambda(L+\mu)\\
	&= 16(L-\lambda)(\mu-\lambda)\\
	&\leq 0
    \end{aligned}
\end{equation}

As $\Delta \leq 0$, the roots of equation \ref{eq:cara} are $a\pm ib$, with $a, b \in \mathbb{R} $.
\begin{itemize}
    \item their module $\rho$ verifies $\rho^2 = a^2+b^2$
    \item their arguments are $ y$ and $-y\in \mathbb{R} $. 
\end{itemize}

The sequence $(a_t)_{t\in \mathbb{R} n}$ is of the form

\begin{equation}
    \label{eq:seq}
    \rho^t(A \cos (ty)+B\sin (ty))
\end{equation}

with $A$ and $B$ real constants. The convergence rate of $a_t$ is thus fully determined by $\rho$. It is thus sufficient to show that $\rho<1$. We know that

\begin{equation}
b= \frac{1}{2} \sqrt{- \Delta}
\end{equation}

Hence

\begin{equation}
    b^2 = -\frac{1}{4} \Delta
\end{equation}

We also have that

\begin{equation}
    a= \frac{1}{2} (1- \gamma\lambda + \beta)
\end{equation}

Hence

\begin{equation}
    a^2 = \frac{1}{4} (1- \gamma\lambda + \beta)^2
\end{equation}

We also have that

\begin{equation}
    \Delta = (1- \gamma\lambda + \beta)^2-4\beta
\end{equation}

Finally

\begin{equation}
    \begin{aligned}
        \label{eq:}
	\rho &=  \sqrt{a^2+b^2}\\
	&= \sqrt{\frac{1}{4} \Big( (1- \gamma\lambda +\beta)^2-\Delta \Big)}\\
	&=  \sqrt{\beta}
    \end{aligned}
\end{equation}

\exercice{}

We have that 
\begin{equation}
	\begin{aligned}
		\label{eq:}
		\sqrt{\beta}^{2t} &\leq (1- \frac{1}{\sqrt{\kappa}})^{2t}\\
		&\leq \exp(- \frac{2t}{\sqrt{\kappa}})
	\end{aligned}
\end{equation}

By a decomposition in a basis of eigenvectors of $H$ (which exists as $H$ is symmetric and real), we obtain that $\theta_t-\theta^*$ converges to $0$ at least at an exponential rate of characteristic time $\sqrt{\kappa}$.


\end{document}
